\documentclass[11pt]{combiArticle}
\usepackage{tikz,amsmath,amsfonts,latexsym}
\usepackage{amssymb}
\usetikzlibrary{shapes,arrows,calc}
\usepackage{calc}
\usepackage{parskip}
\usepackage{enumitem}
\setlist{itemsep=2pt, topsep=2pt, parsep=0pt}
\setlength{\emergencystretch}{2em}

\definecolor{myblue}{RGB}{0, 84, 159}

\lang{de}

\begin{document}
\course{Schichtplan Manager}
\assistent{Linus Palm}
\contact{palm@combi.rwth-aachen.de}
\title{Modell -- Kurzfassung}
\date{09. November 2026}
\maketitle

Der Schichtplan Manager erstellt für einen Monat und eine Station automatisch einen Dienstplan. In der Demo ist
das die erfundene Musterstation. Diese Kurzfassung erklärt ohne Formeln, wie der Plan entsteht. Die genaue
mathematische Beschreibung steht im Dokument \emph{Modell}.

\textbf{Der Ablauf auf einen Blick.}
\begin{enumerate}
  \item \textbf{Eingaben sammeln}: Team, Sollzeiten, Schichten, Mindestbesetzung, Einschränkungen, Wünsche und
        optional der Dienstplan des Vormonats (Abschnitt~\ref{sec:eingaben}).
  \item \textbf{Rechnen}, fest \textbf{30 Sekunden}: Der Solver hält die festen Regeln immer ein
        (Abschnitt~\ref{sec:regeln}) und verbessert die Ziele in einer festen Rangfolge
        (Abschnitt~\ref{sec:ziele}). Ist die Zeit um, gilt der beste bis dahin gefundene Plan.
  \item \textbf{Unabhängig prüfen}: Ein eigenes Prüfprogramm kontrolliert den fertigen Plan noch einmal gegen die
        gesetzlichen und betrieblichen Regeln (Abschnitt~\ref{sec:ergebnis}).
  \item \textbf{Ergebnis anzeigen}: Dienstplan, Besetzung, Lücken, Monatskonten und Kennzahlen.
\end{enumerate}

%==========================================================================================================
\section{Eingaben}\label{sec:eingaben}

\begin{itemize}
  \item \textbf{Monat}, für den geplant wird; immer der ganze Monat.
  \item \textbf{Mitarbeitende} mit ihrer Qualifikation: neun Fachkräfte (sechs in Vollzeit, drei in Teilzeit mit
        80, 85 und 90\,\%), eine Hilfskraft und zwei Azubis, beide in Vollzeit.
  \item \textbf{Sollzeit und Gutschriften} je Person: 38,5 Stunden je Woche, in Teilzeit anteilig, verteilt auf die
        Werktage des Monats ohne Feiertage. Urlaub, Fortbildung und Krankheit werden mit dem Tagessoll
        gutgeschrieben.
  \item \textbf{Schichten} mit Uhrzeiten und den Stunden, die auf das Monatskonto zählen (ohne Pausen):

        \begin{tabular}{l l l}
          \hline
          Schicht & Uhrzeit & zählt \\
          \hline
          Früh (F) & 05:55--13:25 & 7:00 Std. \\
          Zwischen (Z) & 08:30--14:15 & 5:45 Std. \\
          Spät (S) & 13:15--21:00 & 7:15 Std. \\
          Nacht (N) & 20:10--06:10 & 9:15 Std. \\
          \hline
        \end{tabular}
  \item \textbf{Mindestbesetzung} je Tag, Schicht und Qualifikation, jeweils 0 bis 5 Personen. Jeder Monat
        beginnt mit dieser Vorgabe (Personen je Schicht F\,/\,Z\,/\,S\,/\,N):

        \begin{tabular}{l l l}
          \hline
          Qualifikation & Montag bis Freitag & Wochenende und Feiertag \\
          \hline
          Fachkraft & 3\,/\,0\,/\,2\,/\,1 & 2\,/\,0\,/\,1\,/\,1 \\
          Hilfskraft & 1\,/\,0\,/\,0\,/\,0 & -- \\
          Azubi & 1\,/\,0\,/\,1\,/\,0 & -- \\
          \hline
        \end{tabular}
  \item \textbf{Einschränkungen}: Urlaub, Fortbildung, krank, freier Tag, „nur bestimmte Schichten“.
  \item \textbf{Wünsche}: freier Tag, bestimmte Schicht frei, Wunschtag oder Wunschschicht.
  \item Optional der gespeicherte \textbf{Dienstplan des Vormonats}. Er muss unverändert, genau der Vormonat und
        angenommen sein. Mit ihm gelten die Regeln auch über den Monatswechsel; ohne ihn gelten die Tage vor dem
        Monat als frei.
\end{itemize}

Früh- und Spätdienste kann jede Qualifikation übernehmen, auch über ihre Mindestbesetzung hinaus. Nacht- und
Zwischendienste arbeitet eine Qualifikation nur, wenn sie dafür im Monat irgendwann Bedarf hat. Mit der Vorgabe
haben Hilfskraft und Azubis also keine Nacht- und Zwischendienste.

%==========================================================================================================
\section{Feste Regeln}\label{sec:regeln}

Jeder erstellte Dienstplan hält diese Regeln ein:
\begin{itemize}
  \item Höchstens \textbf{ein Dienst pro Person und Tag}; der Dienst zählt mit der Qualifikation der Person.
  \item \textbf{Einschränkungen} gelten immer. Vor einem Abwesenheitstag gibt es keinen Nachtdienst, weil die
        Nacht in ihn hineinreicht. Am Tag vor Urlaubsbeginn und vor einer Fortbildung gibt es keinen Spätdienst:
        Die Fortbildung ist Arbeitszeit, und ihr Beginn könnte sonst in die 11 Stunden Ruhezeit fallen.
  \item Höchstens \textbf{10 Arbeitstage in Folge}; Fortbildungstage zählen mit.
  \item \textbf{Nachtdienste} kommen in Blöcken von 2 bis 3 Nächten. Danach folgen \textbf{2 Tage ganz frei}, also
        auch ohne Fortbildung, und am dritten Tag kein Frühdienst. Zwischen zwei Nachtblöcken liegen mindestens
        7 Tage.
  \item \textbf{Wochenenden}: Samstag und Sonntag dieselbe Schicht oder beide frei; nie zwei Wochenenden in Folge,
        auch nicht über den Monatswechsel. Ein Wochenende gilt als gearbeitet, wenn am Samstag oder Sonntag ein
        Dienst beginnt; ein Nachtdienst am Freitag allein zählt nicht.
  \item \textbf{Endet ein Monat an einem Samstag}, wird der Sonntag danach unsichtbar mitgeplant: als Kopie des
        Samstags, für den alle Regeln gelten. So bekommt am Samstag nur Dienst, wer am Sonntag weitermachen könnte,
        also z.\,B.\ keine dritte Nacht am Samstag. Der Sonntag zählt nicht fürs Monatskonto und steht nicht im Plan;
        geplant wird er im Folgemonat.
  \item Dienste des Vormonats werden \textbf{nicht verändert}, zählen aber bei den Regeln mit.
\end{itemize}

Drei Dinge sind bewusst \emph{keine} festen Regeln, sondern Ziele, damit auch bei Personalmangel immer ein Plan
entsteht:
\begin{itemize}
  \item die \textbf{Mindestbesetzung}: Was fehlt, wird als \emph{Lücke} ausgewiesen;
  \item die \textbf{Sollzeit} und das \textbf{Wochenende über den Monatswechsel}: Sie sind die obersten Ziele und
        werden eingehalten, wann immer das möglich ist;
  \item der \textbf{Wechsel vom Spät- in den Frühdienst}: erlaubt, aber so selten wie möglich.
\end{itemize}

%==========================================================================================================
\section{Ziele in ihrer Rangfolge}\label{sec:ziele}

Die Ziele werden nacheinander verbessert, das wichtigste zuerst. Ein späteres Ziel darf ein früheres nicht
verschlechtern. Nur bei den drei Verteilungszielen (Lücken verteilen, gleichmäßig besetzte Tage, gleichmäßig
besetzte Schichten) ist ein kleiner Spielraum erlaubt, damit die folgenden Ziele noch Wahlmöglichkeiten haben.
\begin{enumerate}
  \item \textbf{Sollzeit einhalten}: Am Monatsende liegt jede Person höchstens \textbf{5 Stunden} über oder unter
        ihrem Soll. Weil dieses Ziel ganz oben steht, hält jeder Plan es ein, sobald das für alle erreichbar ist;
        es wirkt dann wie eine feste Regel. Nur wenn es unmöglich ist, etwa weil jemand den halben Monat frei
        hat (freie Tage werden nicht gutgeschrieben), entsteht trotzdem ein Plan mit möglichst kleiner Abweichung, und die Prüfung weist sie aus.
  \item \textbf{Wochenende über den Monatswechsel fortführen}: Beginnt der Monat an einem Sonntag und ist der
        Vormonat geladen, bekommt der Sonntag die Schicht vom Samstag des Vormonats; wer dort frei hatte, hat auch
        Sonntag frei. Das gilt wie eine feste Regel, sobald es möglich ist. Nur wenn nicht, etwa weil nach dem Plan
        des Vormonats noch Urlaub am Ersten eingetragen wurde, bleibt das Wochenende geteilt, und die Prüfung weist es
        aus.
  \item \textbf{Möglichst wenige Lücken}, also unbesetzte Pflichtstellen.
  \item \textbf{Lücken gleichmäßig verteilen}: Bleiben Lücken, fallen sie gemessen am Bedarf gleichmäßig auf die
        Tage, statt sich an einzelnen Tagen zu häufen.
  \item \textbf{Gleichmäßig besetzte Tage}: Jeder Tag ist im gleichen Verhältnis zu seiner Mindestbesetzung besetzt.
  \item \textbf{Samstage gerecht verteilen}: Alle mit derselben Qualifikation arbeiten möglichst gleich viele
        Samstage, und damit Wochenenden; die Zahlen unterscheiden sich höchstens um eins, soweit Urlaub und
        Einschränkungen es zulassen.
  \item \textbf{Wünsche erfüllen}.
  \item \textbf{Abwechslung}: In jedem Zeitraum von 28 Tagen hat jede Person möglichst jede Schichtart, die sie
        arbeiten kann (mit der Vorgabe: Fachkräfte Früh, Spät und Nacht, Hilfskraft und Azubis Früh und Spät).
        Mit Vormonat reichen diese Zeiträume in den Vormonat zurück; was dort schon gearbeitet wurde, zählt mit.
  \item \textbf{Gleichmäßig besetzte Schichten}: Zusätzliches Personal eines Tages verteilt sich im Verhältnis
        der Mindestbesetzung auf die Schichten.
  \item \textbf{Möglichst wenige Wechsel vom Spät- in den Frühdienst}: Zwischen 21:00 und 05:55 liegen nur knapp
        9 Stunden.
  \item \textbf{Möglichst oft dieselbe Schicht wie am Vortag}, damit ruhige Folgen entstehen.
\end{enumerate}

%==========================================================================================================
\section{Prüfung und Ergebnis}\label{sec:ergebnis}

\textbf{Unabhängige Prüfung.} Ein eigenes Prüfprogramm, das vom Solver nichts übernimmt, kontrolliert den
fertigen Plan gegen gesetzliche und betriebliche Regeln, unter anderem: Mindestbesetzung und Lücken,
Einschränkungen samt keiner Nacht vor einer Abwesenheit und keinem Spätdienst vor Urlaub oder Fortbildung, ein
Dienst pro Tag, Monatssaldo innerhalb von $\pm$\,5 Stunden, tägliche Höchstarbeitszeit und
Pausen, 11 Stunden Ruhezeit zwischen zwei Diensten, höchstens 3 Nächte in Folge und die 48 Stunden Erholung
danach, höchstens 10 Arbeitstage am Stück, am Wochenende dieselbe Schicht an beiden Tagen und nie zwei Wochenenden
in Folge, auch über den Monatswechsel, wenn der Vormonat geladen ist. Den Wechsel vom Spät- in den Frühdienst lässt
die Demo-Version zu und zählt ihn nur. Das Ergebnis ist \emph{angenommen},
\emph{abgelehnt} (mit Person, Datum und Regel) oder \emph{unvollständig geprüft}. Nur ein angenommener Plan kann
als Vormonat für den nächsten Monat geladen werden.

\textbf{Ausgabe.}
\begin{itemize}
  \item \textbf{Der Dienstplan}: für jede Person und jeden Tag die Schicht.
  \item \textbf{Besetzung und Lücken}: je Tag und Schicht die Besetzung im Vergleich zur Mindestbesetzung
        (Ist/Soll) und welche Pflichtstellen mit welcher Qualifikation offen bleiben.
  \item \textbf{Monatskonten}: Soll, Gutschriften, geplante Stunden und Saldo je Person.
  \item \textbf{Kennzahlen}: Lücken, wie viele Monatssalden außerhalb von $\pm$\,5 Stunden liegen und die größte
        Abweichung, erfüllte Wünsche und Wechsel vom Spät- in den Frühdienst.
  \item \textbf{Wünsche}: jeder als \emph{erfüllt}, \emph{nicht erfüllt} oder \emph{nicht erfüllbar}, wenn ihn
        keine Planung erfüllen könnte (z.\,B.\ eine Nacht kurz vor einer Fortbildung oder ein Spätdienst am Tag
        davor).
\end{itemize}
Der Plan lässt sich als \textbf{eine Datei speichern} („Dienstplan \textless Monat\textgreater{}
\textless Jahr\textgreater{} - Version \textless n\textgreater.json“) und später wieder laden, auch als Vormonat
für den nächsten Monat.

\end{document}
